\clearpage
\fancyhead[L]{\small DURCHMESSER DER SONNE.}
\hypertarget{native-524-lower}{}\label{native-524-lower}
\begin{center}\small Göttingische gelehrte Anzeigen. Stück 113. S. 1121..1126. 17. Juli 1809.\\[3pt]\rule{.72\linewidth}{.3pt}\end{center}

In einem Schreiben vom 8. Juni an unseren Herrn Professor \textsc{Gauss} hat Herr \textsc{von Lindenau} die \emph{Ergebnisse von Untersuchungen über den Durchmesser der Sonne} mitgeteilt, die besondere Aufmerksamkeit verdienen. Zu diesen Untersuchungen kam er, als er versuchte, mithilfe einer recht beträchtlichen Anzahl von Sonnenkulminationen, die an der Seeberger Sternwarte beobachtet worden waren, den Sonnendurchmesser aus der Dauer der Durchgänge zu bestimmen. Dabei fand er Unterschiede, die er nur mit Bedenken allein Beobachtungsfehlern zugeschrieben hätte, da sie sehr deutlich einer bestimmten \emph{Periode} zu folgen schienen. Er entschloss sich daher, zu diesem Zweck sämtliche Greenwich-Beobachtungen von 1750 bis 1786 auszuwerten. Die Ergebnisse dieser Untersuchung sind in der Tat sehr bemerkenswert. Herr \textsc{von Lindenau} fasste zunächst für die einzelnen Jahre die Ergebnisse für die Sonnendurchmesser zusammen, die durchgehend auf die mittlere Entfernung der Erde von der Sonne reduziert wurden, wie die folgende Tabelle zeigt:

\SourceLeaf{525}\begin{center}\small
\setlength{\tabcolsep}{1.75pt}
\begin{tabular}{c|c|c||c|c|c}
\shortstack{Jahr der\\Beobach-\\tung}&\shortstack{Mittleres Ergebnis\\für den Sonnen-\\halbmesser bei\\mittlerer\\Entfernung}&\shortstack{Anzahl der\\Beobach-\\tungen}&\shortstack{Jahr der\\Beobach-\\tung}&\shortstack{Mittleres Ergebnis\\für den Sonnen-\\halbmesser bei\\mittlerer\\Entfernung}&\shortstack{Anzahl der\\Beobach-\\tungen}\\ \hline
1750 & $962''63$ & 37 & 1773 & $961''42$ & 56 \\
1751 & $961.82$ & 42 & 1774 & $961.22$ & 32 \\
1752 & $961.61$ & 61 & 1775 & $960.86$ & 62 \\
1753 & $961.81$ & 53 & 1776 & $960.75$ & 53 \\
1754 & $962.30$ & 71 & 1777 & $960.45$ & 81 \\
1755 & $961.81$ & 46 & 1778 & $960.36$ & 71 \\
1765 & $962.44$ & 65 & 1779 & $960.92$ & 41 \\
1766 & $962.70$ & 46 & 1780 & $960.80$ & 37 \\
1767 & $961.57$ & 77 & 1781 & $960.09$ & 50 \\
1768 & $961.53$ & 72 & 1782 & $960.26$ & 37 \\
1769 & $961.23$ & 65 & 1783 & $959.84$ & 42 \\
1770 & $961.66$ & 47 & 1785 & $959.93$ & 65 \\
1771 & $961.82$ & 64 & 1786 & $959.65$ & 55 \\
1772 & $961.84$ & 56 & & & \\
\end{tabular}\end{center}

Es ist sehr auffällig, dass der Sonnenhalbmesser hier stufenweise immer weiter abnimmt. Eine so beträchtliche Abnahme des Sonnenkörpers selbst in so kurzer Zeit ist jedoch völlig undenkbar. Bei einer Abnahme des Sonnenhalbmessers um 3 Sekunden würde sie fast ein Hundertstel des gesamten Volumens ausmachen. Den Grund dieses Phänomen kann man daher wohl nur im Fernrohr, im Auge des Beobachters oder in der Beobachtungsweise suchen. — Ebenso auffällige Ergebnisse lieferte die Zusammenstellung der einzelnen Bestimmungen nach den verschiedenen Jahreszeiten. Während dieser ganzen Reihe von Jahren zeigten sich mit nur wenigen Ausnahmen unverkennbare, periodisch wiederkehrende Änderungen. Die mittleren Werte für die einzelnen Monate, durchweg auf den 15. reduziert, waren folgende:

\begin{center}\begin{tabular}{l|c||l|c}
&\shortstack{Halbmesser\\der Sonne}&&\shortstack{Halbmesser\\der Sonne}\\ \hline
Januar & $960''17$ & Juli & $960''14$ \\
Februar & $961.16$ & August & $961.00$ \\
März & $961.52$ & September & $961.70$ \\
April & $961.22$ & Oktober & $961.80$ \\
Mai & $961.20$ & November & $961.16$ \\
Juni & $960.00$ & Dezember & $960.43$ \\
\end{tabular}\end{center}

Es zeigten sich hier zwei Maxima, die man auf Anfang April und Anfang Oktober legen kann, und zwei Minima, die man auf Anfang Januar und Anfang Juli legen kann. Der Unterschied zwischen den Extremwerten lässt sich für den Halbmesser auf anderthalb und damit für den Durchmesser auf drei Sekunden veranschlagen. Manchem mag es freilich etwas bedenklich erscheinen, über so kleine Größen anhand von \emph{Zeitbeobachtungen} entscheiden zu wollen, die für sich genommen oft viel größere Unterschiede zeigen. Andererseits sprechen die große Genauigkeit, mit der \textsc{Bradley} und \textsc{Maskelyne} in Greenwich immer beobachtet haben, die außerordentlich große Zahl der Beobachtungen und die Regelmäßigkeit der Änderungen in den Ergebnissen zu sehr dagegen, jene Unterschiede für bloß zufällig und für Folgen von Beobachtungsfehlern zu halten. Man könnte zunächst daran denken, sie dem Einfluss der Jahreszeiten auf die Atmosphäre und damit einer Veränderung der Irradiation oder einem Einfluss auf die Empfindlichkeit des Auges zuzuschreiben. Dann müssten die beiden verschiedenen Extremwerte allem Anschein nach aber in die Zeiten der größten Wärme und Kälte fallen, während die Beobach\SourceLeaf{518}tungen für beide Zeiten denselben Durchmesser ergeben. Durch eine \emph{elliptische Gestalt des Sonnenkörpers} ließe sich das Phänomen dagegen sehr ungezwungen erklären. Man nehme an, die Sonne sei ein durch Drehung um ihre Rotationsachse erzeugtes Ellipsoid, und bezeichne den Halbmesser des Sonnenäquators mit $a$ und die halbe Achse mit $b$. Ferner bezeichne man die Neigung des Sonnenäquators gegen den Erdäquator mit $i$, die Rektaszension des aufsteigenden Knotens der ersten Ebene auf der zweiten mit $N$ und die Rektaszension der Sonne mit $A$. Schließlich sei der scheinbare Halbmesser einer Kugel mit dem Halbmesser $a$, aus der Entfernung $1$ gesehen, mit $r$ bezeichnet. Eine einfache Herleitung zeigt dann, dass der aus der beobachteten Kulminationszeit abgeleitete und auf die mittlere Entfernung reduzierte Halbmesser folgender ist:

\[ r\left(1-\tfrac{1}{2}\left(1-\frac{bb}{aa}\right)\cos(A-N)^2\sin i^2\right) \]

Aus der von früheren Astronomen bestimmten Lage des Sonnenäquators, dessen aufsteigender Knoten auf der Ekliptik bei einer Länge von $78^\circ$ und dessen Neigung gegen die Ekliptik mit $7\tfrac{1}{2}$ Grad angenommen wird, folgt $i = 26^\circ \, 2'$, $N = 16^\circ \, 55'$. Die äußersten Werte des Halbmessers treten also ein, wenn $A-N$ entweder $0$ oder $90^\circ$ oder $180^\circ$ oder $270^\circ$ ist. Im ersten und dritten Fall, also am 8. April und 12. Oktober, ist jener Halbmesser $= r\left(1-\left(1-\dfrac{bb}{aa}\right).\,0{,}09631\right)$. Im zweiten und vierten Fall dagegen, also am 8. Juli und 6. Januar, findet man den Halbmesser $= r$. Der erste Wert ist der kleinere oder der größere, je nachdem, ob die Sonne ein abgeplattetes oder ein längliches Sphäroid ist. Sehr sonderbar ist, dass die beobachteten Phänomene bei dieser Erklärung die letztere Hypothese erfordern. Tatsächlich lassen sie sich sehr gut darstellen, wenn man $\dfrac{bb}{aa}-1 = \tfrac{1}{58}$ oder $\dfrac{b}{a} = \tfrac{17}{116}$ setzt oder für die Äquatorialabplattung des Sonnenkörpers $\tfrac{1}{117}$ ansetzt. Herr \textsc{von Lindenau} findet etwas weniger, nämlich $\tfrac{1}{140}$. — Um zu sehen, was direkte Messungen über diese Gestalt entscheiden würden, leitete Herr \textsc{von Lindenau} nun auch die vertikalen Sonnendurchmesser aus den Unterschieden der Zenitdistanzen der beiden Ränder ab. Dazu benutzte er ebenso viele Beobachtungen. Die Ergebnisse sind folgende:

\begin{center}\begin{tabular}{c|c||c|c}
Jahr & \shortstack{Mittlerer\\Sonnenhalbmesser} & Jahr & \shortstack{Mittlerer\\Sonnenhalbmesser} \\ \hline
1765 & $963''81$ & 1775 & $963''15$ \\
1766 & $964.28$ & 1776 & $962.71$ \\
1767 & $963.88$ & 1777 & $961.74$ \\
1768 & $963.40$ & 1778 & $962.47$ \\
1769 & $964.07$ & 1779 & $962.85$ \\
1770 & $963.31$ & 1780 & $962.63$ \\
1771 & $962.89$ & 1781 & $961.94$ \\
1772 & $963.36$ & 1782 & $961.84$ \\
1773 & $963.16$ & 1783 & $961.52$ \\
1774 & $963.46$ & & \\
\end{tabular}\end{center}

Auch hier zeigt sich wie bei den horizontalen Durchmessern eine schrittweise Abnahme. Die vertikalen Durchmesser sind durchweg um einige Sekunden größer als die horizontalen, im Mittel um $3''65$. Herr \textsc{von Lindenau} berechnete daraus eine Äquatorialabplattung von $\tfrac{1}{288}$. Herr Professor \textsc{Gauss} findet noch etwas weniger, nämlich $\tfrac{1}{420}$. Sowohl der Vergleich der zu verschiedenen Jahreszeiten gemessenen horizontalen Sonnendurchmesser untereinander als auch der Vergleich der mittleren horizontalen mit den mittleren vertikalen Durchmessern sprechen also übereinstimmend für eine längliche Gestalt der Sonne, sofern man sich auf die Genauigkeit der Messungen selbst verlassen kann, obwohl die Werte für die Elliptizität nicht übereinstimmen. Je weniger man versteht, wie bei einer solchen Gestalt des Sonnenkörpers ein Gleichgewichtszustand an seiner Oberfläche möglich sein soll, umso gespannter wird \SourceLeaf{519}man darauf sein, bald die Einzelheiten dieser interessanten Untersuchung von Herrn \textsc{von Lindenau} kennenzulernen.

Wir nutzen noch die Gelegenheit, ein Verzeichnis mitzuteilen, das Herr \textsc{von Wisniewski}, außerordentliches Mitglied der Petersburger Akademie der Wissenschaften, aus Astracan an Herrn Dr. \textsc{Olbers} gesandt hat und das uns dieser übermittelt hat. Das Verzeichnis nennt einige Sternbedeckungen durch den Mond, die noch in diesem Jahr eintreten und in \textsc{Bode}s Astronomischem Jahrbuch fehlen. Wir wollen die Beobachter darauf aufmerksam machen und so vielleicht einige Beobachtungen gewinnen, die denen entsprechen, die Herr \textsc{von Wisniewski} im Russischen Reich anstellen wird. Das Verzeichnis ist für den Horizont eines Ortes auf der Breite $52^{\circ}\,31'$ und der östlichen Länge $2^{\mathrm u}\,0^{\mathrm m}$ von Greenwich berechnet. — —

Mehrere dieser Bedeckungen sind bereits in der Connoissance des tems angezeigt.
